\refitem{ASSOCIAÇÃO BRASILEIRA DE NORMAS TÉCNICAS. \textbf{NBR 6024}: informação
e documentação: numeração progressiva das seções de um documento: apresentação.
Rio de Janeiro: ABNT, 2012.}

\refitem{ASSOCIAÇÃO BRASILEIRA DE NORMAS TÉCNICAS. \textbf{NBR 14724}:
informação e documentação: trabalhos acadêmicos: apresentação. Rio de Janeiro:
ABNT, 2011.}

\refitem{BRASIL. Tribunal de Contas da União. \textbf{Criando documentos digitais
acessíveis}. Brasília: TCU, 2010. Disponível em:
\url{https://www.tjdft.jus.br/acessibilidade/publicacoes/producao_de_conteudo_com_acessibilidade_vf3.pdf}.
Acesso em: 14 dez. 2021.}

\refitem{NASCIMENTO, Maria Vanessa do.; MARTINS, Graci Kelly. A trajetória das
escolas de biblioteconomia no Brasil. \textbf{REBECIN}: Revista Brasileira de
Educação em Ciência da Informação, São Paulo, v.4, n. esp., p. 37-54, 2. sem.,
2017. Disponível em: \url{https://portal.abecin.org.br/rebecin/article/view/90}.
Acesso em: 14 dez. 2021.}

\refitem{SALTON, Bruna Poletto; AGNOL, Anderson Dall; TUCARTTI, Alissa.
\textbf{Manual de acessibilidade em documentos digitais}. Bento Gonçalves:
Instituto Federal de Educação, Ciência e Tecnologia do Rio Grande do Sul, 2017.}

\refnota{O autor deve descrever todas as referências que foram citadas dentro da monografia.
As referências devem ser alinhadas à esquerda e com espaçamento simples. São
organizadas em ordem alfabética por sobrenome dos autores. Entre cada
referência, deve ter uma linha em branco com espaço simples, de acordo com a
NBR 6023/2018.}
